\documentclass[12pt]{article}
\usepackage[utf8]{inputenc}
\usepackage[T1]{fontenc}
\usepackage{hyperref}
\title{Port Sandwich Language Resources for K-12 Curriculum Development}
\author{Compiled by Claude Sonnet 5.5 (claude-sonnet-5-5) for Andr\'as Kornai}
\date{}
\begin{document}
\maketitle

\section{Introduction}
Port Sandwich (ISO 639-3 \texttt{psw}, Glottocode port1285, also called Lamap) is an Oceanic language of south-east Malekula, Vanuatu. Sources found give about 1{,}000--1{,}200 speakers. The main documentation is by Jean-Michel Charpentier, in French. No school materials were found.

\section{Primary Educational Resources}
None found.

\section{Linguistic and Reference Materials}
\begin{itemize}
\item J.-M. Charpentier (1979), \emph{La Langue de Port-Sandwich (Nouvelles-H\'ebrides): Introduction phonologique et grammaire}, CNRS, Paris, 208 pp. (catalogue entry). \url{https://glottolog.org/resource/reference/id/131706}
\item J.-M. Charpentier (1974), \emph{Langue de Port-Sandwich}, dissertation, Universit\'e de Bordeaux III (catalogue entry). \url{https://glottolog.org/resource/reference/id/134295}
\item Wikipedia, ``Port Sandwich language''. \url{https://en.wikipedia.org/wiki/Port_Sandwich_language}
\item Grambank feature data for Port Sandwich. \url{https://grambank.clld.org/languages/port1285}
\end{itemize}

\section{Religious and Community Resources}
None found.

\section{Digital and Audio Resources}
The Endangered Languages Project page for Port Sandwich was returned; its content (audio or otherwise) was not verified. \url{https://endangeredlanguages.com/node/30115}

\section{Curriculum Development Context}
Documentation is largely limited to Charpentier's French-language grammar and a 1974 Port-Sandwich--French dictionary mentioned in references (no online text found). A curriculum would require community-led development, building from these descriptions.

\end{document}